
% --- quiz1 ---
\element{essai}{
  \begin{questionmult}{q:quiz1}
What are the signal characteristics and the initial objectives described in the $\text{Lab}$ on digital filtering?
    \begin{choices}
      \correctchoice{The studied signal is corrupted by a slow drift ($\text{drift}$) of its mean value, a common problem in sensor measurements.}
      %answer: The signal is corrupted by a slow drift, which is a common problem with sensor measurements.
      %tip: Check the initial description of the signal.
      \correctchoice{One of the objectives of the $\text{Lab}$ is to design a filter to remove this slow drift.}
      %answer: A first filter will be designed in order to remove this drift.
      %tip: The first filtering objective is the removal of the drift.
      \wrongchoice{The sampling frequency specified for the signal analysis is $\text{Fs}=800 \text{ Hz}$.}
      %answer: The specified sampling frequency is $\text{Fs}=8000 \text{ Hz}$.
      %tip: Check the numerical value of the sampling frequency.
      \wrongchoice{The signal is periodic with a period of about $40 \text{ Hz}$ in the frequency domain.}
      %answer: The signal is periodic with a period of about 40 in the time domain, and not in frequencies.
      %tip: The period is a time measurement, not a frequency measurement.
    \end{choices}
  \end{questionmult}
}

% --- quiz2 ---
\element{nocategory}{
  \begin{questionmult}{q:quiz2}
According to the time and frequency analysis of signal $\text{x}$, which observations are correct?
    \begin{choices}
      \correctchoice{The time signal $\text{x}$ is approximately periodic with a period of about $40$ samples.}
      %answer: We observe an approximately periodic time series, with a period of about 40.
      %tip: The analysis of the signal in the time domain is described.
      \correctchoice{The Fourier Transform ($\text{FT}$) of the signal consists of a series of regularly spaced spectral lines, indicating a periodic underlying signal.}
      %answer: We see that the Fourier Transform consists of a series of regularly spaced spectral lines, which is an indication that the underlying signal is periodic.
      %tip: The appearance of the $\text{FT}$ is related to the periodicity of the signal.
      \wrongchoice{The $\text{FT}$ of the signal shows a spectral peak spacing of $0.025$ in non-normalized frequencies (Hertz).}
      %answer: The peaks are spaced by $0.025$ in normalized frequencies.
      %tip: The spacing is given in normalized frequencies.
      \wrongchoice{The analysis of the $\text{FT}$ shows only the slow drift component without any other spectral lines.}
      %answer: The $\text{FT}$ shows a series of regularly spaced spectral lines and a slow drift (average value).
      %tip: The signal contains several frequency components.
    \end{choices}
  \end{questionmult}
}

% --- quiz54 (pythontex template) ---
\element{nocategory}{
  \begin{questionmultx}{q:quiz54}
  \pyc{a = rng.normal(0, 1)}
  \pyc{b = rng.normal(0, 1)}
  \pyc{expected_0 = round(a+b, 4)}
  \pyc{_nbch_0 = len(str(abs(int(expected_0))))}
  \pyc{_ndec_0 = len(str(expected_0).split(".")[-1].rstrip("0")) if "." in str(expected_0) else 0}
  \pyc{_ndigits_0 = _nbch_0 + _ndec_0}
  \pyc{expected_1 = round(a-b, 4)}
  \pyc{_nbch_1 = len(str(abs(int(expected_1))))}
  \pyc{_ndec_1 = len(str(expected_1).split(".")[-1].rstrip("0")) if "." in str(expected_1) else 0}
  \pyc{_ndigits_1 = _nbch_1 + _ndec_1}
  \py{%
  rf"This is a numeric question where you must calculate the sum and difference of {a:.4f} and {b:.4f}"
     }
  {%
  \def\AMCbeginQuestion#1#2{}%
  \AMCquestionNumberfalse
    \begin{questionmultx}{q:quiz54_0}
Sum: 
      \pys{\AMCnumericChoices{!{expected_0}}{digits=!{_ndigits_0},decimals=!{_ndec_0}}}
  \end{questionmultx}
  }
  {%
  \def\AMCbeginQuestion#1#2{}%
  \AMCquestionNumberfalse
    \begin{questionmultx}{q:quiz54_1}
Difference: 
      \pys{\AMCnumericChoices{!{expected_1}}{digits=!{_ndigits_1},decimals=!{_ndec_1}}}
  \end{questionmultx}
  }
  \end{questionmultx}
}

% --- quiz59 ---
\element{nocategory}{
  \begin{questionmultx}{q:quiz59}
Please enter the mean and standard deviation values of the time series above
  {%
  \def\AMCbeginQuestion#1#2{}%
  \AMCquestionNumberfalse
    \begin{questionmultx}{q:quiz59_0}
Mean: 
  \AMCnumericChoices{0.0}{digits=2,decimals=1}
  \end{questionmultx}
  }
  {%
  \def\AMCbeginQuestion#1#2{}%
  \AMCquestionNumberfalse
    \begin{questionmultx}{q:quiz59_1}
Standard deviation 
  \AMCnumericChoices{1.0}{digits=2,decimals=1}
  \end{questionmultx}
  }
  \end{questionmultx}
}

% --- quiz63 (pythontex template) ---
\element{nocategory}{
  \begin{questionmult}{q:quiz63}
  \pyc{n = rng.integers(1, 10)}
  \pyc{v_exp_0 = bool(True)}
  \pyc{v_exp_1 = bool(False)}
  \pyc{v_exp_2 = bool(False)}
  \pyc{v_exp_3 = bool(False)}
  \py{%
  rf"What is the derivative of $x^n$, with n={n}, that is  $x^{{{n}}}$,"
     }
    \begin{choices}
\py{%
    rf"\correctchoice{{ The derivative is ${n} x^{{{n-1}}}$ }}"
    if v_exp_0
    else
    rf"\wrongchoice{{ The derivative is ${n} x^{{{n-1}}}$ }}"
    }
      %answer: 
      %tip: 
\py{%
    rf"\correctchoice{{ The derivative is $\frac{{{1}}}{{{n}}} x^{{{n}}}$ }}"
    if v_exp_1
    else
    rf"\wrongchoice{{ The derivative is $\frac{{{1}}}{{{n}}} x^{{{n}}}$ }}"
    }
      %answer: 
      %tip: 
\py{%
    rf"\correctchoice{{ The derivative is $\frac 1{{{n-1}}} x^{{{n-1}}}$ }}"
    if v_exp_2
    else
    rf"\wrongchoice{{ The derivative is $\frac 1{{{n-1}}} x^{{{n-1}}}$ }}"
    }
      %answer: 
      %tip: 
\py{%
    rf"\correctchoice{{ The derivative is ${n-1} x^{{{n-1}}}$ }}"
    if v_exp_3
    else
    rf"\wrongchoice{{ The derivative is ${n-1} x^{{{n-1}}}$ }}"
    }
      %answer: 
      %tip: 
    \end{choices}
  \end{questionmult}
}

% --- quiz79 (pythontex template) ---
\element{test}{
  \begin{questionmultx}{q:quiz79}
  \pyc{FN = rng.integers(0, 10)}
  \pyc{FP = rng.integers(0, 10)}
  \pyc{TN = rng.integers(0, 50)}
  \pyc{TP = rng.integers(0, 50)}
  \pyc{expected_0 = round(TP/(TP+FN), 3)}
  \pyc{_nbch_0 = len(str(abs(int(expected_0))))}
  \pyc{_ndec_0 = len(str(expected_0).split(".")[-1].rstrip("0")) if "." in str(expected_0) else 0}
  \pyc{_ndigits_0 = _nbch_0 + _ndec_0}
  \pyc{expected_1 = round((TP+TN)/(TP+TN+FP+FN), 3)}
  \pyc{_nbch_1 = len(str(abs(int(expected_1))))}
  \pyc{_ndec_1 = len(str(expected_1).split(".")[-1].rstrip("0")) if "." in str(expected_1) else 0}
  \pyc{_ndigits_1 = _nbch_1 + _ndec_1}
  \py{%
  rf"Consider the following \emph{{confusion matrix}}:
{{\def\LTcaptype{{none}} % do not increment counter
\begin{{longtable}}{{|l|l|l|}}
\toprule\noalign{{}}
True/Predict & 0 & 1 \\
\midrule\noalign{{}}
\endhead
\bottomrule\noalign{{}}
\endlastfoot
0 & {TN} & {FP} \\
1 & {FN} & {TP} \\
\end{{longtable}}
}}"
     }
  {%
  \def\AMCbeginQuestion#1#2{}%
  \AMCquestionNumberfalse
    \begin{questionmultx}{q:quiz79_0}
Give the recall (tolerance $\pm$ 0.01)
      \pys{\AMCnumericChoices{!{expected_0}}{digits=!{_ndigits_0},decimals=!{_ndec_0}}}
  \end{questionmultx}
  }
  {%
  \def\AMCbeginQuestion#1#2{}%
  \AMCquestionNumberfalse
    \begin{questionmultx}{q:quiz79_1}
Give the accuracy (tolerance $\pm$ 0.01)
      \pys{\AMCnumericChoices{!{expected_1}}{digits=!{_ndigits_1},decimals=!{_ndec_1}}}
  \end{questionmultx}
  }
  \end{questionmultx}
}

% --- quiz81 (pythontex template) ---
\element{test}{
  \begin{questionmultx}{q:quiz81}
  \pyc{w0 = rng.normal(0, 1)}
  \pyc{w1 = rng.normal(0, 1)}
  \pyc{x = rng.normal(0, 1)}
  \pyc{expected = round(round(abs(w1),3)*round(x,3) + round(w0,3),3)}
  \pyc{_nbch = len(str(abs(int(expected))))}
  \pyc{_ndec = len(str(expected).split(".")[-1].rstrip("0")) if "." in str(expected) else 0}
  \pyc{_ndigits = _nbch + _ndec}
  \py{%
  rf"Consider a scatter plot such as the one shown in the figure
\begin{{center}}
\includegraphics[width=0.38\linewidth,height=\textheight,keepaspectratio,alt={{Scatter plot}}]{{./Figure_1.png}}
\\[-0.6em]
{{\small Figure:~ \itshape Scatter plot}}
\end{{center}}
We perform a linear regression \texttt{{reg}}, which yields - \texttt{{reg.intercept\_}} = {w0:.3f} and - \texttt{{reg.coef\_}} = {abs(w1):.3f}. For a new observation x = {x:.3f}, what would be the value predicted by the model?"
     }
Give the accuracy (tolerance $\pm$ 0.01)
    \pys{\AMCnumericChoices{!{expected}}{digits=!{_ndigits},decimals=!{_ndec}}}
  \end{questionmultx}
}

% --- quiz82 (pythontex template) ---
\element{test}{
  \begin{questionmult}{q:quiz82}
  \pyc{a = rng.normal(0, 1)}
  \pyc{b = rng.normal(0, 1)}
  \pyc{x = rng.normal(0, 1)}
  \pyc{x2 = rng.normal(0, 1)}
  \pyc{v_exp_0 = bool(False)}
  \pyc{v_exp_1 = bool(True)}
  \pyc{v_exp_2 = bool(a > 0)}
  \py{%
  rf"Consider the simple regression  model 
$$
y = \alpha x + \beta
$$
with $\alpha={{{round(a,3)}}}$ and $\beta={{{round(b,3)}}}$"
     }
    \begin{choices}
\py{%
    rf"\correctchoice{{ For x={round(x2,3)}, the prediction is {round(b,3)*round(x2,3) + round(a,3) :.3f} }}"
    if v_exp_0
    else
    rf"\wrongchoice{{ For x={round(x2,3)}, the prediction is {round(b,3)*round(x2,3) + round(a,3) :.3f} }}"
    }
      %answer: Prediction is simply computed by $y = \alpha x + \beta$ 
      %tip: 
\py{%
    rf"\correctchoice{{ For x={round(x,3)}, the prediction is {round(a,3)*round(x,3) + round(b,3) :.3f}  
 }}"
    if v_exp_1
    else
    rf"\wrongchoice{{ For x={round(x,3)}, the prediction is {round(a,3)*round(x,3) + round(b,3) :.3f}  
 }}"
    }
      %answer: Prediction is simply computed by $y = \alpha x + \beta$
      %tip: 
\py{%
    rf"\correctchoice{{ The slope is positive }}"
    if v_exp_2
    else
    rf"\wrongchoice{{ The slope is positive }}"
    }
      %answer: The slope is actually {'not' if a<0 else ''} positive, since it is equal to $\alpha$={a:.3f}
      %tip: 
    \end{choices}
  \end{questionmult}
}